\newpage
\section{\huge \textcolor{blue}{一阶线性ODE}}
\subsection{\LARGE 一阶线性ODE标准方程}
对于如下方程形式:
{\Large
\begin{equation}
a(x) y' + b(x) y = c(x)
\end{equation}
}
我们称 $y$ 和 $y'$ 呈现出线性关系, 可类比如下:
{\Large
\begin{equation}
a y_1 + b y_2 = c
\end{equation}
}
特别地, 若$c=0$, 则称为齐次方程, 目前暂不考虑\\
下对(10)式进行变形:\\ 
同除$a(x)$, 
并记$b(x)/a(x)=p(x)$, $c(x)/a(x)=q(x)$, 得:
{\Large
\begin{equation}
y' + p(x) y = q(x)
\end{equation}
}
同时能注意到, 部分教程使用如下方程:
{\Large
\begin{equation}
y' = p(x) y + q(x)
\end{equation}
}
此方程表达更突出"一阶", 为后续笔记的统一性
和计算的便利, 不使用此形式

\newpage
\subsection{\LARGE 温度模型}

下文列举一个实际例子(温度传导模型)来说明:
\begin{figure}[htbp]
    \centering
\includegraphics[width=0.6\textwidth]{C:/Users/37/Desktop/Notes on Common Calculus Equations/Notes and pictures/cwffc3.1.1.png}
    \caption{传导模型}
    \label{fig:cwffc3.1.1}
\end{figure}

类似如图的温度传导模型, 我们能给出相关的牛顿定律:
{\Large
\begin{equation}
\frac{dT}{dt} = k (T_e - T)
\end{equation}
}
其中$k > 0$, $t = time$, $T = temperature$, 且有$T_e > T$, $T(0) = T_0$\\
下将其写为标准方程(一般性的线性方程):
{\Large
\begin{equation}
\frac{dT}{dt} + k T = k T_e
\end{equation}
}


\subsection{\LARGE 解一阶线性方程}
为解决上述模型的方程问题, 即需解决下式:
{\Large
\begin{equation}
y' + p(x) y = q(x)
\end{equation}
}
为此介绍如下方法:
\subsubsection{\large 积分因子法}
对方程两边同乘以一个函数$\mu(x)$, 得:
{\Large
\begin{equation}
\mu(x) y' + \mu(x) p(x) y = \mu(x) q(x)
\end{equation}
}
将等式左边变形为乘积求导形式:
{\Large \[ (\mu(x) y)' = \mu(x) q(x) \]}
现要使上式成立, 则需满足:
{\Large \[ \mu'(x) = p(x) \mu(x) \]}
即求:
{\Large \[ \frac{\mu'(x)}{\mu(x)} = p(x) \]}
解得:
{\Large
\begin{equation}
\mu(x) = e^{\int p(x) dx}
\end{equation}
}
现得到$\mu(x)$, 称为积分因子\\
下用一道例题展示该方法:\\
eg 3.1.1: \hspace*{6em} $xy' - y = x^3$\\
\hspace*{2em}解
{\begin{align*}
    & \qquad y' - \frac{1}{x}y = x^2 \quad (x \neq 0) \\
    & \qquad e^{-\ln x} = (e^{\ln x})^{-1} = \frac{1}{x} \\
    & \qquad \therefore \frac{1}{x}y' - \frac{1}{x^2}y = x \\
    & \qquad \left( \frac{1}{x} y \right)' = x \\
    & \qquad \therefore \frac{1}{x} \cdot y = \frac{1}{2}x^2 + C \\
    & \qquad \therefore y = \frac{1}{2}x^3 + Cx
\end{align*}
}
了解积分因子法后, 我们再回到温度传导模型中, 解决$p(x)$为常数的情况:
{\Large \[ \frac{dT}{dt} + k T = k T_e \]}
则有$\mu(x)$为:
{\Large \[ \mu(t) = e^{\int k dt} = e^{kt} \]}
最后解得:
{\Large
\begin{equation}
T = e^{-kt} \int k T_e e^{kt} dt + C e^{-kt}
\end{equation}
}
此解称为温度传导的暂态解, 可用于预测物体温度随时间变化的情况, 
当$t \to \infty$, 即物体温度趋近于环境温度时, 得稳态解:
{\Large
\begin{equation}
T = e^{-kt} \int k T_e e^{kt} dt
\end{equation}
}

